\clearpage
\section{Appendix: Scope and Qualifications}
This concise study treats the Twenty-ninth Sunday in Ordinary Time, Year A, 18 October 2026, in the dioceses of the United States: Roman Missal, Third Edition, U.S. English (2011), based on the third typical Latin edition, emended reprint (2008); U.S. Lectionary for Mass, second typical edition, volume I, no.145. No local-calendar overlay was specified; Saint Luke supplies no added proper.

All ten ritual positions and both Communion alternatives are identified. The Missal-antiphon and Lectionary routes are studied; permitted chant substitutions, sprinkling, the coupled Preface/Eucharistic Prayer choices and a concluding solemn blessing remain unselected. The Sunday Gloria and Creed belong to the Ordinary. No Sequence, proper Offertory chant or uniquely appointed Preface is supplied. Exact rubrical branches appear in the expansive study's Appendix: Liturgical Resolution.

The whole-formulary arguments are editorial syntheses over the Fathers' and saints' identified passages. Ordinary Time's apostolic reading follows an independent semicontinuous course; the shared Week XXIX prayers are not peculiar to Year A. The expansive historical appendix named on page 2 supplies the chronology evidence and its qualifications.

Scriptural English quoted in the discussion follows the public-domain Douay--Rheims/Challoner study witness, not the approved proclaimed English; the U.S. Lectionary controls proclamation. Approved Missal and Lectionary bodies are neither reproduced nor reconstructed. Incipits and analytical descriptions are not substitute prayers. Exact 2008 Latin and publisher-specific 2011 altar-book pages remain uninspected; the shared owner's layered digital, official-excerpt and publisher evidence supports this bounded citation-led treatment, not complete altar-book collation.

Jerome's DCO text has no established historical exemplar; Corpus Thomisticum supplies electronic witnesses rather than inspected historical printings. Bede's decisive passage was collated with PL 92, cols.236A--B; his complete digital commentary is not thereby critically collated. NPNF translations retain their editorial layers and Augustine's abridgment. The bounded reception sweep makes no claim to exhaustive scholarship. This work is a study aid, not an official liturgical book or ecclesiastical approval.

\section{References}
Augustine's \work{Expositions} are cited by Vulgate Psalm numbers XVI, XXXII and XCV; the historical NPNF I.8 English headings are XVII, XXXIII and XCVI. Its CCEL delivery gives an 1886 print basis while its preface is dated 1888. Chrysostom's Matthew homilies use Prevost's translation revised by M.B. Riddle, NPNF I.10 (1888); his Pauline homilies use the Oxford translation revised for NPNF I.13 (1889).

Jerome: \work{Commentarii in Isaiam}, XII, DCO Latin \work{Pars 3}, pp.126--129. Aquinas: Corpus Thomisticum Latin expositions at the paragraph numbers cited; \work{Summa} III, q.48, a.4 in the English Dominican translation, Gutenberg 19950 (the electronic ``Art.3'' below ``Fourth Article'' is a heading error). Bede: \work{In Marci}, III, PL 92 (Paris, 1862), cols.236A--B. Gregory: Browne--Swallow, NPNF II.7 (1894). John Paul II: \href{https://www.vatican.va/content/john-paul-ii/en/audiences/2002/documents/hf_jp-ii_aud_20020918.html}{General Audience, 18 September 2002}, §§1--5, original paraphrase.

Appointment and occurrence: USCCB \work{Liturgical Calendar 2026}, pp.5,41, and dated readings for 18 October; \work{Ordo lectionum Missae} (1981), no.145 A; General Introduction to the Lectionary §§65--69,103--104; U.S. GIRM §§48,51,53,61--63,68,74,87--90,364--365. Official controls checked 13 October 2026. The expansive historical appendix identifies the Catholic Encyclopedia, NABRE and 1912 PBC chronology witnesses by article and locus.
